% Source-derived chapter 06: Compatibility with class field theory
% Original source: tates-thesis-de-rham-setting
\chapter{Compatibility with class field theory}
\label{tates-thesis-de-rham-setting-chapter-06}
\source{tates-thesis-de-rham-setting}

\begin{remark}[Source section]
\label{tates-thesis-de-rham-setting-chapter-06-source}
\source{tates-thesis-de-rham-setting}
\source{tates-thesis-de-rham-setting}
This chapter reproduces the corresponding section of the retrieved
LaTeX source, including its definitions, statements, examples, and
proofs. It has not yet been translated into Lean.
\notready
\end{remark}

% The source section title was: Compatibility with class field theory
\label{s:cft}
\source{tates-thesis-de-rham-setting}

In this section, we show that the construction of \S \ref{s:weyl} is compatible
with geometric class field theory in a suitable sense. 

\subsection{Conventions regarding $D$-modules}\label{ss:diff-1}
\source{tates-thesis-de-rham-setting}

\subsubsection{} Before proceeding, we take a moment to establish some
notational conventions regarding $D$-modules. 
We refer to \cite{crystals} for details.

\subsubsection{}

Let $X$ be a smooth variety. As in \emph{loc. cit}., 
we have the so-called \emph{left} forgetful functor 
$\Oblv^{\ell}:D(X) \to \QCoh(X)$. It is normalized
so that the functor $\Oblv^{\ell}[\dim X]$
is $t$-exact; for example, this functor sends the dualizing
$\omega_X$ to $\sO_X$ (with its standard left $D$-module structure).

\subsubsection{}

The functor $\Oblv^{\ell}$ admits the left adjoint
$\ind^{\ell}$. For us, we take as a definition that the sheaf of
differential operators on $X$ is:
%
\[
D_X \coloneqq \ind^{\ell}(\sO_X).
\] 

\noindent We remark that this object is concentrated in degree $-\dim X$.
(On general grounds, it coincides with $\ind^r(\omega_X)$, where 
$\ind^r$ is the right $D$-module induction functor and
$\omega_X = \Omega_X^{\dim X}[\dim X] \in \QCoh(X)$ is the
dualizing sheaf.)

\subsubsection{}

Now suppose $X$ is smooth and affine. Let $\on{Diff}(X)$ be the 
algebra of global differential operators on $X$.

As in \cite{crystals}, $\End_{D(X)}(D_X)$ canonically coincides with
$\on{Diff}(X)$, but with reversed multiplication. It follows that there
is a canonical equivalence:
%
\begin{equation}\label{eq:diff-dmod}
\source{tates-thesis-de-rham-setting}
\on{Diff}(X)\mod \simeq D(X)
\end{equation}

\noindent between left modules for $\on{Diff}(X)$ and the category $D(X)$,
sending $\on{Diff}(X)$ to $D_X$. We remark that this functor is only
$t$-exact up to shift. 

\subsection{Construction of functors}

\subsubsection{}

In \S \ref{s:weyl}, we constructed compact objects
$\sF_n \in \IndCoh^*(\sY^{\leq n})^c$ and morphisms:
%
\[
W_n^{op} \to \ul{\End}_{\IndCoh^*(\sY^{\leq n})}(\sF_n).
\]

\noindent We obtain a corresponding functor:
%
\[
\ol{\Delta}_n:W_n\mod 
\overset{\eqref{eq:diff-dmod}}{\simeq} 
D(\fL_n^+\bA^1) \to \IndCoh^*(\sY_n)
\in \DGCat_{cont}
\]

\noindent uniquely characterized by sending $D_{\fL_n^+\bA^1}$ to 
$\sF_n$ compatibly with the action of $W_n^{op} = 
\TwoEnd_{D(\fL_n^+\bA^1)}(D_{\fL_n^+\bA^1})$. 

\begin{rem}
\source{tates-thesis-de-rham-setting}
\notready

Above, we carefully normalized $D_{\fL_n^+ \bA^1}$ to lie in cohomological
degree $-n$. This is not especially relevant in our analysis 
until \S \ref{s:main}.

\end{rem}


\subsubsection{}\label{sss:delta-n}
\source{tates-thesis-de-rham-setting}

For later use, we also use the notation 
$\Delta_n$ to denote the further composition:
%
\[
D(\fL_n^+\bA^1) \xar{\ol{\Delta}_n} \IndCoh^*(\sY_n) 
\xar{\lambda_{n,*}^{\IndCoh}} \IndCoh^*(\sY).
\]

\subsection{Cartier duality}

We now review some constructions from geometric class field theory.

\subsubsection{}

Let $\sK_n \subset \fL^+ \bG_m$ be the $n$th congruence subgroup, 
i.e., $\sK_n \coloneqq \Ker(\fL^+\bG_m \to \fL_n^+\bG_m)$.

Recall that there is a canonical bimultiplicative pairing:
%
\begin{equation}\label{eq:cartier}
\source{tates-thesis-de-rham-setting}
\LocSys_{\bG_m}^{\leq n} \times (\fL\bG_m/\sK_n)_{dR} \to \bB \bG_m.
\end{equation}

\noindent One basic property is that its restriction along the map:
%
\[
\LocSys_{\bG_m}^{\leq n} \times \fL\bG_m/\sK_n
\to 
\LocSys_{\bG_m}^{\leq n} \times (\fL\bG_m/\sK_n)_{dR} \to \bB \bG_m
\]

\noindent factors through the Contou-Carr\`ere duality pairing:
%
\[
\LocSys_{\bG_m}^{\leq n} \times \fL\bG_m/\sK_n \to 
\bB \fL^{\leq n} \bG_m \times \fL\bG_m/\sK_n \xar{CC} 
\bB \bG_m.
\]

\noindent Here we normalize the sign in the Contou-Carr\`ere pairing
so that we have a commutative diagram:
%
\[
\xymatrix{
\bB \fL \bG_m \times \fL^+\bG_m \ar[rr] \ar[d] & & 
\bB \fL \bG_m \times \fL\bG_m \ar[d]^{CC} \\
\bB \Gr_{\bG_m} \times \fL^+\bG_m \ar[rr] & &
\bB \bG_m
}
\]

\noindent where the bottom arrow is induced by the pairing of 
Proposition \ref{p:cc-log}.

From \eqref{eq:cartier} and functoriality, we obtain a symmetric
monoidal Fourier-Mukai functor:
%
\begin{equation}\label{eq:trun-lcft}
\source{tates-thesis-de-rham-setting}
D(\fL\bG_m/\sK_n) \to \QCoh(\LocSys_{\bG_m}^{\leq n}).
\end{equation}

\begin{thm}
\source{tates-thesis-de-rham-setting}
\notready\label{t:trun-lcft}
\source{tates-thesis-de-rham-setting}

The functor \eqref{eq:trun-lcft} is an equivalence.
In particular, there is a canonical fully faithful symmetric monoidal functor
$D(\fL_n^+\bG_m) \into \QCoh(\LocSys_{\bG_m}^{\leq n})$.

\end{thm}

Passing to the limit over $n$, we also obtain the following theorem
of Beilinson-Drinfeld.

\begin{thm}
\source{tates-thesis-de-rham-setting}
\notready\label{t:full-lcft}
\source{tates-thesis-de-rham-setting}

There is a canonical symmetric monoidal equivalence:
%
\[
D^*(\fL\bG_m) \isom \QCoh(\LocSys_{\bG_m}).
\]

\end{thm}

\subsection{Formulation of the equivariance property}

Now observe that $\ol{\Delta}_n$ maps a category acted on by 
$D(\fL_n^+\bG_m)$ to one acted on by 
$\QCoh(\LocSys_{\bG_m}^{\leq n})$ (via the structure map 
$\sY^{\leq n} \to \LocSys_{\bG_m}^{\leq n}$).

Using Theorem \ref{t:trun-lcft}, we can regard both sides as 
acted on by $D(\fL_n^+\bG_m)$. 

In the remainder of this section, we show that 
$\ol{\Delta}_n$ is canonically a morphism of 
$D(\fL_n^+\bG_m)$-module categories.

\subsection{Digression on Harish-Chandra data}\label{ss:hc}
\source{tates-thesis-de-rham-setting}

We will construct the equivariance structure using the theory of \emph{Harish-Chandra
data}. We review this theory below.

In what follows, $G$ is an affine algebraic group and $A \in \Alg$ is a DG algebra.
In our applications, $A$ will be classical, so at some points in the discussion we 
will assume that.

\subsubsection{} 

First, suppose that $G$ acts on $A$ as an associative algebra. In other words,
we assume we are given a lift of $A$ along the forgetful functor 
$\Alg(\Rep(G)) \to \Alg(\Vect) = \Alg$.

This defines a canonical weak action of $G$ on $A\mod$. Its basic property is that
the forgetful functor $A\mod \to \Vect$ is weakly $G$-equivariant.

\subsubsection{}

Now suppose that $\sC$ is a DG category with a weak $G$-action. 

Recall that $\Rep(G)$ naturally acts on $\sC^{G,w}$.
Therefore, we can consider $\sC^{G,w}$ as enriched over $\Rep(G)$.
For $\sF,\sG\in \sC^{G,w}$, we let:
%
\[
\ul{\Hom}_{\sC}^{enh}(\sF,\sG) \in \Rep(G)
\]

\noindent denote the corresponding mapping complex.

\begin{rem}
\source{tates-thesis-de-rham-setting}
\notready

There are natural maps:
%
\[
\Oblv\ul{\Hom}_{\sC}^{enh}(\sF,\sG) \to \ul{\Hom}_{\sC}(\Oblv(\sF),\Oblv(\sG))
\]
%
\noindent that are easily seen to be isomorphisms for $\sF$ compact. 

\end{rem}

Similarly, for
$\sF \in \sC^{G,w}$, we let:
%
\[
\ul{\End}_{\sC}^{enh}(\sF,\sG) \in \Alg(\Rep(G)) 
\]

\noindent denote the inner endomorphism algebra.

\subsubsection{}\label{sss:weak-functor}
\source{tates-thesis-de-rham-setting}

In our earlier setting, suppose we are given a weakly $G$-equivariant
functor:
%
\[
F:A\mod \to \sC.
\]

\noindent Note that $A \in A\mod^{G,w}$, so we obtain a natural object 
$\sF \coloneqq F^{G,w}(A) \in \sC^{G,w}$. Moreover, there is a canonical map:
%
\[
\vph:A \simeq \ul{\End}_{A\mod}^{enh}(A) \to \ul{\End}_{\sC^{G,w}}(\sF)^{op} \in \Alg(\Rep(G)). 
\]

\noindent (Here the superscript \emph{op} indicates the reversed multiplication.) 

Conversely, given $\sF \in \sC^{G,w}$ and $\vph$ as above, we obtain a functor 
$F$ as above. Indeed, as is standard, a datum $\vph$ is the
same as specifying a morphism 
%
\[
A\mod^{G,w} = A\mod(\Rep(G)) \to \sC^{G,w} 
\]

\noindent of $\Rep(G)$-module categories; by de-equivariantization
(i.e., tensoring over $\Rep(G)$ with $\Vect$), we obtain the 
desired functor $F:A\mod \to \sC$.

\subsubsection{}\label{sss:hc-recap}
\source{tates-thesis-de-rham-setting}

Next, recall from \cite{methods} \S 10
that $\Alg(\Rep(G))$ carries a canonical monad $B \mapsto U(\fg) \# B$.

Here $U(\fg) \# B$ is the usual smash (or crossed) product construction.
Non-derivedly, modules for $U(\fg) \# B$ are vector spaces equipped
with an action of $B$ and an action of $\fg$ that are compatible under the
action of $\fg$ on $B$ by derivations. For a proper derived construction 
(in the more complicated setting of topological algebras), we refer to
\cite{methods}.

By definition, making $A \in \Alg(\Rep(G))$ into a module over this monad
is a \emph{Harish-Chandra} datum for $A$. When $A$ is classical, this 
amounts to the data of a morphism:
%
\[
i:\fg \to A 
\]

\noindent such that $i$ is a $G$-equivariant morphism of Lie algebras such that 
$[i(\xi),-]:A \to A$
coincides with the derivation defined by  $\xi \in \fg$ and the $G$-action on $A$.

From \emph{loc. cit}., it follows that specifying a Harish-Chandra datum for
$A$ is the same as extending the weak $G$-action on $A\mod$ to a strong $G$-action.

\subsubsection{}\label{sss:hc-functor}
\source{tates-thesis-de-rham-setting}

Now suppose that $G$ acts strongly on $\sC$. Let $\sF \in \sC^{G,w}$
be a fixed object. We claim that $\ul{\End}_{\sC}^{enh}(\sF) \in \Alg(\Rep(G))$ 
carries a canonical Harish-Chandra datum.

The construction is formal. Let $\sD_0 \subset \sC^{G,w}$ be the
(non-cocomplete) DG category generated by $\sF$ under finite colimits and 
direct summands. As in \cite{dario-*/!}, the monoidal category:
%
\[
\sH\sC_G \coloneqq \TwoEnd_{G\mod}(D(G)^{G,w})^{op} = \fg\mod^{G,w}
\]

\noindent of Harish-Chandra bimodules acts canonically on $\sC^{G,w}$. 
Moreover, $\sH\sC_G$ is \emph{rigid} monoidal, so its (non-cocomplete) 
monoidal subcategory 
$\sH\sC_G^c$ of compact objects preserves $\sD_0$. Therefore, 
$\sH\sC_G$ acts on $\sD_1 \coloneqq \Ind(\sD_0)$. By de-equivariantization,
$\sD_1 = \sD^{G,w}$ for a canonical strong $G$-category $\sD$.
But clearly $\sD = \ul{\End}_{\sC}^{enh}(\sF)\mod$ by construction,
so we obtain the Harish-Chandra datum as desired (from \S \ref{sss:hc-recap}).

\begin{rem}
\source{tates-thesis-de-rham-setting}
\notready

In particular, there is a canonical map 
$\fg \to \ul{\End}_{\sC}^{enh}(\sF)$ of Lie algebras in $\Rep(G)$.
This map is the standard \emph{obstruction to strong equivariance}
on $\sF$.

\end{rem}

Moreover, we see that if $A \in \Alg(\Rep(G))$ is equipped with a Harish-Chandra
datum, giving a strongly $G$-equivariant functor:
%
\[
F:A\mod \to \sC
\]

\noindent is equivalent to giving $\sF \in \sC^{G,w}$ and
a map $\vph:A \to \ul{\End}_{\sC}^{enh}(\sF)^{op}$ that is
\emph{compatible with the Harish-Chandra data}.

In the case that $A$ is classical and $\ul{\End}_{\sC}^{enh}(\sF)^{op}$ 
is coconnective, this simply amounts to saying that composition:
%
\[
\fg \xar{i} A \xar{\vph} \ul{\End}_{\sC}^{enh}(\sF)^{op}
\]

\noindent coincides with the obstruction to strong equivariance on the object $\sF$.

\subsubsection{}\label{sss:connected-obstr}
\source{tates-thesis-de-rham-setting}

We now add one more mild observation before concluding our discussion. The
reader may safely skip this material and return back to it as necessary.

Suppose $G$ acts weakly on $\sC$ and $\sF,\sG \in \sC^{G,w}$.
There is a canonical map:

\[
\Oblv(\ul{\Hom}_{\sC}^{enh}(\sF,\sG)) \to \ul{\Hom}_{\sC}(\Oblv(\sF),\Oblv(\sG))
\in \Vect
\]

\noindent that is an isomorphism whenever $\sF$ is compact. (This is the
reason we use the subscript $\sC$ in $\ul{\Hom}_{\sC}^{enh}$ and
not $\ul{\Hom}_{\sC^{G,w}}^{enh}$: the underlying vector space
of the representation $\ul{\Hom}_{\sC}^{enh}(\sF,\sG)$ is 
comparable to, and often isomorphic to, the complex of maps
from $\sF$ to $\sG$ in $\sC$ itself, not in $\sC^{G,w}$.)

Now suppose that we are in the setting of \S \ref{sss:hc-functor}.
Suppose $\sF \in \sC^{G,w}$ is compact, $A$ is classical, and
$\ul{\End}_{\sC}^{enh}(\sF)$ is coconnective. Then the data
of a map:
%
\[
A \to \ul{\End}_{\sC}^{enh}(\sF)^{op}
\]

\noindent compatible with the Harish-Chandra data is equivalent to 
a map:
% 
\[
\vph:A \to  \ul{\End}_{\sC}(\sF)^{op} \in \Alg = \Alg(\Vect)
\]

\noindent such that the composition:
%
\[
\fg \xar{i} A \xar{\vph} \ul{\End}_{\sC}(\sF)^{op}
\]

\noindent is the obstruction to strong equivariance. Indeed, this follows
from the assumption that $\sF$ is compact, so $\ul{\End}_{\sC}(\sF)^{op} = 
\ul{\End}_{\sC}^{enh}(\sF)^{op}$, and the connectedness of $G$, so that
the map $\vph$ will automatically be a morphism of $G$-representations.
(We remark that the co/connectivity assumptions allow us to replace 
$\ul{\End}_{\sC}(\sF)^{op}$ with the classical algebra 
$H^0\ul{\End}_{\sC}(\sF)^{op}$.)
 
\subsection{Reduction}

We now return to our particular setting, and apply the above material to 
reduce our construction to a calculation. 

\subsubsection{}

Let $G = \fL_n^+\bG_m$ act strongly on $\sC \coloneqq \IndCoh^*(\sY^{\leq n})$
via geometric class field theory. By construction, we have:
%
\[
\IndCoh^*(\sY^{\leq n})^{\fL_n^+\bG_m,w} \simeq  
\IndCoh^*(\sY_{\log}^{\leq n}).
\]

\noindent Under this dictionary, the forgetful functor 
$\sC^{G,w} \to \sC$ corresponds to the $\IndCoh$-pushforward 
$\sY_{\log}^{\leq n} \to \sY^{\leq n}$ (i.e., $!$-averaging
for $\Gr_{\bG_m}^{\leq n}$).

\subsubsection{}

We have the object $i_{n,*}(\sO_{\sZ^{\leq n}}) \in \IndCoh^*(\sY_{\log}^{\leq n})$
that maps to $\sF_n \in \IndCoh^*(\sY^{\leq n})$ under this forgetful functor.
In \S \ref{s:weyl}, we constructed the map:
%
\[
\vph:A \coloneqq W_n \to \ul{\End}_{\sC}^{enh}(\sF_n)^{op} \simeq 
\ul{\End}_{\sC}(\sF_n)^{op}.
\]

\noindent (We remind that the displayed isomorphism is a formal consequence
of compactness of $\sF_n$.)

Now on the one hand, we have a morphism: 
%
\[
\Lie(\fL_n^+\bG_m) \simeq k[[t]]/t^n \to \Gamma(\fL_n^+\bA^1,T_{\fL_n^+\bA^1}) 
\to W_n
\]

\noindent encoding the infinitesimal action of $\fL_n^+\bG_m$ on
$\fL_n^+\bA^1$. We denote this morphism by $\alpha$.

On the other hand, we have the composition:
%
\[
\Lie(\fL_n^+\bG_m) \simeq k[[t]]/t^n \simeq 
(t^{-n} k[[t]]dt/k[[t]] dt)^{\vee} \subset 
\Gamma(\fL^{pol,\leq n} \bA^1 dt,\sO_{\fL^{pol,\leq n} \bA^1 dt})
\to \ul{\End}_{\sC}^{enh}(\sF_n)^{op}
\]

\noindent where the last morphism here uses the map:
%
\[
\sZ^{\leq n} \subset \sY_{\log}^{\leq n} \to \LocSys_{\bG_m,\log}^{\leq n}
\xar{\Pol} \fL^{pol,\leq n} \bA^1 dt.
\]

\noindent We denote this morphism by $\beta$.

In \S \ref{ss:hc-calc-pf}, we will prove:

\begin{lem}
\source{tates-thesis-de-rham-setting}
\notready\label{l:hc-calc}
\source{tates-thesis-de-rham-setting}

The diagram:
%
\[
\xymatrix{
\Lie(\fL_n^+\bG_m) \simeq k[[t]]/t^n \ar[d]^{\alpha} \ar[drr]^{\beta} \\
W_n \ar[rr]^{\vph} && \ul{\End}_{\sC}(\sF_n)
}
\]

\noindent commutes.

\end{lem}

Observe that under Theorem \ref{t:trun-lcft}, $\beta$ is interpreted
as the obstruction to strong equivariance. Therefore, 
from \S \ref{sss:connected-obstr}, we deduce:

\begin{prop}
\source{tates-thesis-de-rham-setting}
\notready\label{p:eq}
\source{tates-thesis-de-rham-setting}

The functor $\ol{\Delta}_n$ is canonically a morphism of
$D(\fL_n^+\bG_m)$-module categories. 

\end{prop}

\subsection{Proof of Lemma \ref{l:hc-calc}}\label{ss:hc-calc-pf}
\source{tates-thesis-de-rham-setting}

\subsubsection{}

Using the additive structure on $\fL_n^+\bA^1$,
we compute its global vector fields as:
%
\[
\Gamma(\fL_n^+\bA^1,T_{\fL_n^+\bA^1}) \simeq 
k[[t]]/t^n \otimes \Sym(t^{-n}k[[t]]dt/k[[t]]dt).
\]

\noindent A straightforward calculation shows that the infinitesimal action
map:
%
\[
\Lie(\fL_n^+\bG_m) \simeq k[[t]]/t^n \to \Gamma(\fL_n^+\bA^1,T_{\fL_n^+\bA^1})
\simeq k[[t]]/t^n \otimes \Sym(t^{-n}k[[t]]dt/k[[t]]dt)
\]

\noindent which was used in the definition of $\alpha$ above,
is the map:
%
\[
\begin{gathered}
(f \in k[[t]]/t^n) \mapsto (1 \otimes f) \cdot \cotr = (f \otimes 1) \cdot \cotr
\in k[[t]]/t^n \otimes t^{-n}k[[t]]dt/k[[t]]dt \\ \subset 
k[[t]]/t^n \otimes \Sym(t^{-n}k[[t]]dt/k[[t]]dt).
\end{gathered}
\]

\noindent Here we recall the canonical vector:
%
\[
\begin{gathered}
\cotr \in 
\Gamma(\sD^{\leq n},\sO_{\sD^{\leq n}}) \otimes
\Gamma(\sD^{\leq n},\sO_{\sD^{\leq n}})^{\vee} = \\
\Gamma^{\IndCoh}(\sD^{\leq n} \times \sD^{\leq n},
\sO_{\sD^{\leq n}} \boxtimes \omega_{\sD^{\leq n}}) = 
k[[t]]/t^n \otimes t^{-n}k[[t]]dt/k[[t]]dt
\end{gathered}
\]

\noindent from Lemma \ref{l:cotr-1}.

Therefore, we have:
%
\[
\alpha(f) = (1 \otimes f) \cdot \cotr \in 
k[[t]]/t^n \otimes t^{-n}k[[t]]dt/k[[t]]dt \subset W_n.
\]

\subsubsection{}

To simplify the discussion, we take $f = 1$ for the time being. 
We wish to show that $\vph\alpha(1) = \beta(1)$, i.e., we wish 
to show that $\vph(\cotr) = \beta(1)$.
We will adapt the calculation to discuss the
case of general $f$ in \S \ref{sss:gen'l-f}.

\subsubsection{}

Let us unwind the construction of $\vph(\cotr)$ in this case.

First, we form:
%
\[
\cotr \in \Gamma^{\IndCoh}(\sD^{\leq n} \times \sD^{\leq n},
\sO_{\sD^{\leq n}} \boxtimes \omega_{\sD^{\leq n}}).
\]

\noindent Second, we form:\footnote{Here we abuse notation: in \S \ref{ss:act-inv}, we used the
notation $\omega^{univ}$ for a similar construction with $\sZ^{\leq n}$ in place of
$\fL^{pol,\leq n} \bA^1 dt$; that setting is pulled back from the present
one.}
%
\[
\omega^{univ} \in \Gamma^{\IndCoh}(\sD^{\leq n} \times 
\fL^{pol,\leq n} \bA^1 dt,
\omega_{\sD^{\leq n}} \boxtimes \sO_{\fL^{pol,\leq n} \bA^1 dt}).
\]

\noindent We then form:

\[
p_{12}^*(\cotr) \cdot p_{23}^*(\omega^{univ}) \in 
\Gamma^{\IndCoh}(\sD^{\leq n} \times \sD^{\leq n} \times 
\fL^{pol,\leq n} \bA^1 dt,
\omega_{\sD^{\leq n}} \boxtimes \omega_{\sD^{\leq n}} \boxtimes 
\sO_{\fL^{pol,\leq n} \bA^1 dt}).
\]

According to \S \ref{sss:recipe}, $\vph(\cotr)$ is calculated as follows.
We pull back the above section to $\sD^{\leq n} \times \sD^{\leq n} \times 
\sZ^{\leq n}$, pushforward along:
%
\[
\sD^{\leq n} \times \sD^{\leq n} \times \sZ^{\leq n}
\xar{\AJ^{-1} \times \AJ \times \id}
\Gr_{\bG_m}^{\leq n} \times \Gr_{\bG_m}^{\leq n} 
\times \sZ^{\leq n} \xar{\on{mult} \times \id} 
\Gr_{\bG_m}^{\leq n} \times \sZ^{\leq n}
\] 

\noindent and apply the construction of \S \ref{ss:gpd}.

Of course, the pullback and the pushforward commute here. 
Therefore, in what follows we calculate the pushforward with $\sZ^{\leq n}$
replaced by $\fL^{pol,\leq n} \bA^1 dt$.

In what follows, let:
%
\[
\sigma \in \Gamma^{\IndCoh}(\Gr_{\bG_m}^{\leq n} \times \fL^{pol,\leq n} \bA^1 dt,
\omega_{\Gr_{\bG_m}^{\leq n}} \boxtimes \sO_{\fL^{pol,\leq n} \bA^1 dt})
\]

\noindent denote the resulting section. So our objective, in effect, is to calculate
$\sigma$.

\subsubsection{}

Recall from Proposition \ref{p:cc-log} that: 
%
\begin{equation}\label{eq:gr-cc}
\source{tates-thesis-de-rham-setting}
\Gamma^{\IndCoh}(\Gr_{\bG_m}^{\leq n},\omega_{\Gr_{\bG_m}^{\leq n}}) \simeq 
\Gamma(\fL_n^+ \bG_m,\sO_{\fL_n^+ \bG_m}).
\end{equation}

\noindent Moreover, this is an isomorphism of commutative algebras. It is
convenient to reinterpret the above construction using this isomorphism.

By construction, note that the composition:

\[
t^{-n}k[[t]]dt/k[[t]]dt \simeq 
\Gamma^{\IndCoh}(\sD^{\leq n},\omega_{\sD^{\leq n}}) \to 
\Gamma^{\IndCoh}(\Gr_{\bG_m}^{\leq n},\omega_{\Gr_{\bG_m}^{\leq n}}) 
\simeq \Gamma(\fL_n^+ \bG_m,\sO_{\fL_n^+ \bG_m})
\]

\noindent using pushforward along $\AJ$ sends:
%
\[
\eta \in t^{-n}k[[t]]dt/k[[t]]dt
\] 

\noindent to the function:
%
\[
(g \in \fL_n^+\bG_m) \mapsto \Res(g \eta).
\]

\noindent Similarly, if we use $\AJ^{-1}$ instead, we obtain the
function:
%
\[
(g \in \fL_n^+\bG_m) \mapsto \Res(g^{-1} \eta).
\]

\subsubsection{}

Our original section:
%
\[
p_{12}^*(\cotr) \cdot p_{23}^*(\omega^{univ}) \in 
\Gamma^{\IndCoh}(\sD^{\leq n} \times \sD^{\leq n} \times 
\fL^{pol,\leq n} \bA^1 dt,
\omega_{\sD^{\leq n}} \boxtimes \omega_{\sD^{\leq n}} \boxtimes 
\sO_{\fL^{pol,\leq n} \bA^1 dt})
\]

\noindent can be interpreted as a map:
%
\[
\begin{gathered}
\fL^{pol,\leq n} \bA^1 dt = 
\ul{t^{-n}k[[t]]dt/k[[t]]dt} \to 
\ul{t^{-n}k[[t]]dt/k[[t]]dt \otimes t^{-n}k[[t]]dt/k[[t]]dt}
\end{gathered}
\]

\noindent where for a finite-dimensional vector
space $V$, we let $\ul{V}$ denote the scheme $\Spec(\Sym(V^{\vee}))$.
It is immediate to see that the above map is given
by the formula:
%
\[
\eta \mapsto (\eta \otimes 1) \cdot \cotr.
\]

Pushing $p_{12}^*(\cotr) \cdot p_{23}^*(\omega^{univ})$ forward along:
%
\[
\sD^{\leq n} \times \sD^{\leq n} \times \sZ^{\leq n}
\xar{\AJ^{-1} \times \AJ \times \id}
\Gr_{\bG_m}^{\leq n} \times \Gr_{\bG_m}^{\leq n} \times \fL^{pol,\leq n} \bA^1 dt
\]

\noindent and applying \eqref{eq:gr-cc}, we obtain a function on 
%
\[
\fL_n^+\bG_m \times \fL_n^+\bG_m \times \fL^{pol,\leq n} \bA^1 dt.
\]

\noindent We deduce that it is given by the formula:
%
\begin{equation}\label{eq:res-res}
\source{tates-thesis-de-rham-setting}
(g_1,g_2,\eta) \mapsto 
(\Res \otimes \Res)((g_1^{-1} \otimes g_2) \cdot 
(\eta \otimes 1) \cdot \cotr).
\end{equation}

\subsubsection{}

We now wish to calculate \eqref{eq:res-res} more explicitly. 

First, we note the following, which is an immediate calculation.

\begin{lem}
\source{tates-thesis-de-rham-setting}
\notready\label{l:res-cotr}
\source{tates-thesis-de-rham-setting}

For $\eta \in \fL^{pol,\leq n} \bA^1 dt$, we have:
%
\[
(\id \otimes \Res)((\eta \otimes 1) \cdot \cotr) = \eta.
\]

\end{lem}

We then deduce:

\begin{lem}
\source{tates-thesis-de-rham-setting}
\notready\label{l:res-res}
\source{tates-thesis-de-rham-setting}

For $(g,\eta) \in  \fL_n^+\bA^1 \times \fL^{pol,\leq n} \bA^1 dt$, we have:
%
\[
(\Res \otimes \Res)((\eta \otimes g) \cdot \cotr) = \Res(g\eta).
\]

\end{lem}

\begin{proof}

We have:
%
\[
(1 \otimes g) \cdot \cotr = (g \otimes 1) \cdot \cotr.
\]

\noindent Indeed, this encodes the fact that the residue pairing is
$k[[t]]$-equivariant.

Therefore, we have:
%
\[
(\eta \otimes g) \cdot \cotr = (g\eta \otimes 1) \cdot \cotr
\]

\noindent which gives:
%
\[
(\Res \otimes \Res)((\eta \otimes g) \cdot \cotr) = 
\Res (\id \otimes \Res)((g\eta \otimes 1) \cdot \cotr).
\]

\noindent Now the result follows from Lemma \ref{l:res-cotr}.

\end{proof}

In the setting of \eqref{eq:res-res}, we deduce:
%
\begin{equation}\label{eq:res-res-calc}
\source{tates-thesis-de-rham-setting}
(\Res \otimes \Res)((g_1^{-1} \otimes g_2) \cdot 
(\eta \otimes 1) \cdot \cotr) = 
\Res(\frac{g_2}{g_1} \eta).
\end{equation}

Now, our section $\sigma$ is obtained by using a pushforward
along $\Gr_{\bG_m}^{\leq n} \times \Gr_{\bG_m}^{\leq n} \to \Gr_{\bG_m}^{\leq n}$.
On the dual side, this means we restrict our function along: 
%
\[
\fL_n^+ \bG_m \xar{\Delta \times \id}
\fL_n^+\bG_m \times \fL_n^+\bG_m \times \fL^{pol,\leq n} \bA^1 dt.
\]

\noindent By \eqref{eq:res-res-calc}, we see the resulting function is:
%
\[
(g,\eta) \mapsto \Res(\eta).
\]

This clearly coincides with the endomorphism defined by $\beta(1)$, giving the
claim (in the $f=1$ case).

\subsubsection{}\label{sss:gen'l-f}
\source{tates-thesis-de-rham-setting}

As promised, we now treat the case of general $f$.

Briefly, one replaces \eqref{eq:res-res} with:
%
\[
(g_1,g_2,\eta) \mapsto 
(\Res \otimes \Res)((g_1^{-1} \otimes g_2) \cdot 
(\eta \otimes 1) \cdot (1 \otimes f) \cdot \cotr),
\]

\noindent which Lemma \ref{l:res-res} implies is:
%
\[
(g_1,g_2,\eta) \mapsto \Res(\frac{f g_2}{g_1}\eta),
\]

\noindent which for $g_1 = g_2 = g$ is:
%
\[
(g,\eta) \mapsto \Res(f \eta).
\]

\noindent Again, the induced endomorphism clearly coincides with that 
defined by $\beta(f)$, concluding the argument.
